\documentclass[dvipdfmx,a4paper,11pt]{article}
\usepackage[utf8]{inputenc}
%\usepackage[dvipdfmx]{hyperref} %リンクを有効にする
\usepackage{url} %同上
\usepackage{amsmath,amssymb} %もちろん
\usepackage{amsfonts,amsthm,mathtools} %もちろん
\usepackage{braket,physics} %あると便利なやつ
\usepackage{bm} %ラプラシアンで使った
\usepackage[top=30truemm,bottom=30truemm,left=25truemm,right=25truemm]{geometry} %余白設定
\usepackage{latexsym} %ごくたまに必要になる
\renewcommand{\kanjifamilydefault}{\gtdefault}
\usepackage{otf} %宗教上の理由でmin10が嫌いなので


\usepackage[all]{xy}
\usepackage{amsthm,amsmath,amssymb,comment}
\usepackage{amsmath}    % \UTF{00E6}\UTF{0095}°\UTF{00E5}\UTF{00AD}\UTF{00A6}\UTF{00E7}\UTF{0094}¨
\usepackage{amssymb}  
\usepackage{color}
\usepackage{amscd}
\usepackage{amsthm}  
\usepackage{wrapfig}
\usepackage{comment}	
\usepackage{graphicx}
\usepackage{setspace}
\usepackage{pxrubrica}
\usepackage{enumitem}
\usepackage{mathrsfs} 
\usepackage{caption}

\setstretch{1.2}


\newcommand{\R}{\mathbb{R}}
\newcommand{\Z}{\mathbb{Z}}
\newcommand{\Q}{\mathbb{Q}} 
\newcommand{\N}{\mathbb{N}}
\newcommand{\C}{\mathbb{C}} 
\newcommand{\Sin}{\text{Sin}^{-1}} 
\newcommand{\Cos}{\text{Cos}^{-1}} 
\newcommand{\Tan}{\text{Tan}^{-1}} 
\newcommand{\invsin}{\text{Sin}^{-1}} 
\newcommand{\invcos}{\text{Cos}^{-1}} 
\newcommand{\invtan}{\text{Tan}^{-1}} 
\newcommand{\Area}{\text{Area}}
\newcommand{\vol}{\text{Vol}}
\newcommand{\maru}[1]{\raise0.2ex\hbox{\textcircled{\tiny{#1}}}}
\newcommand{\sgn}{{\rm sgn}}
%\newcommand{\rank}{{\rm rank}}



   %当然のようにやる．
\allowdisplaybreaks[4]
   %もちろん．
%\title{第1回. 多変数の連続写像 (岩井雅崇, 2020/10/06)}
%\author{岩井雅崇}
%\date{2020/10/06}
%ここまで今回の記事関係ない
\usepackage{tcolorbox}
\tcbuselibrary{breakable, skins, theorems}

\theoremstyle{definition}
\newtheorem{thm}{定理}
\newtheorem{lem}[thm]{補題}
\newtheorem{prop}[thm]{命題}
\newtheorem{cor}[thm]{系}
\newtheorem{claim}[thm]{主張}
\newtheorem{dfn}[thm]{定義}
\newtheorem{rem}[thm]{注意}
\newtheorem{exa}[thm]{例}
\newtheorem{conj}[thm]{予想}
\newtheorem{prob}[thm]{問題}
\newtheorem{rema}[thm]{補足}

\DeclareMathOperator{\Ric}{Ric}
\DeclareMathOperator{\Vol}{Vol}
 \newcommand{\pdrv}[2]{\frac{\partial #1}{\partial #2}}
 \newcommand{\drv}[2]{\frac{d #1}{d#2}}
  \newcommand{\ppdrv}[3]{\frac{\partial #1}{\partial #2 \partial #3}}


%ここから本文．
\begin{document}
%\maketitle

\newpage
\begin{center}
{\Large 2022年度秋冬学期 大阪大学 理学部数学科 \\ 幾何学1演義} \\
金曜4限(15:10-16:40) 理学部E310
\end{center}
\begin{flushright}
 岩井雅崇(いわいまさたか) 2022/10/07 \\
\end{flushright}
{\large 基本的事項}
\begin{itemize}
  \setlength{\parskip}{0cm} % 段落間
  \setlength{\itemsep}{0cm} % 項目間
\item この授業は対面授業です. 金曜4限(15:10-16:40) 理学部E310にて演習の授業を行います.
\item 授業ホームページ(\url{https://masataka123.github.io/2022_winter_stokes/})にて授業の問題等をアップロードしていきます. 
下に授業ホームページのQRコードを貼っておきます. 
\begin{figure}[htbp]
\begin{center}
 \includegraphics[height=30mm, width=30mm]{stokes.png}
 \caption*{授業のQRコード}
\end{center}
\end{figure}
\end{itemize}

\hspace{-18pt}{\large 成績に関して}

次の1と2を満たしているものに単位を与えます.
\begin{enumerate}
  \setlength{\parskip}{0cm} % 段落間
  \setlength{\itemsep}{0cm} % 項目間
\item 幾何学1の講義の単位が可以上である.
\item 14回授業時までに演習の授業で少なくとも2回以上発表している.
\end{enumerate}
ただし2の条件を達成できないものには別途救済レポートを課して2を達成したものとすることがある.

成績は講義の点数に演習の出来具合を加算してつける(と思います).

\vspace{11pt}
\hspace{-18pt}{\large 解答の仕方について}
\begin{itemize}
  \setlength{\parskip}{0cm} % 段落間
  \setlength{\itemsep}{0cm} % 項目間
  \item 問題の解答を黒板に書いて発表してください. 正答だった場合その問題はそれ以降解答できなくなります. 不正解だった場合他の人に解答権が移ります. 
  \item 複数人が解答したい問題があるときは平和的な手段で解答者を決めます. 
  \item 英語問題を答える際には英語を和訳してください. なお解答は日本語で行っても良い.
  \item 演習問題の難易度は一定ではない. 難しい問題と英語問題を解いた場合は成績に加点を行う. 
 \end{itemize}


\vspace{11pt}
\hspace{-18pt}{\large その他}
\begin{itemize}
  \setlength{\parskip}{0cm} % 段落間
  \setlength{\itemsep}{0cm} % 項目間
    \item 授業内容をあまり把握していないので, 演習問題と授業内容が噛み合ってない可能性があります.
  \item 休講あるいは補講をすることがあるので, こまめにホームページとCLEは確認してください.
    \item オフィスアワーを月曜15:00-17:00に設けています. この時間に私の研究室に来ても構いません(ただし来る場合は前もって連絡してくれると助かります.)
 \end{itemize}
 
 

 \end{document}
